\chapter{数字经济与就业、收入分配}

\begin{examguide}
\textbf{本章考情}：数字经济对就业与分配的影响是 0258 论述题的高频主题（\important{专硕·核心}），技能偏向型技术进步（SBTC）的推导是学硕宏观与劳动经济学交叉考点（\important{学硕·热点}）；数字鸿沟、平台用工是政策热点素材。

\textbf{核心考点}：（1）替代效应与创造效应的分解；（2）SBTC 模型与技能溢价的推导；（3）任务模型与自动化边界；（4）就业极化与工资极化的辨析；（5）技能需求逆转的证据与机制；（6）零工经济与平台用工关系；（7）超级明星效应、劳动收入份额与数据租金；（8）数字鸿沟三层结构。

\textbf{本章主线}：技术如何影响就业（9.1）→ 就业结构如何极化（9.2）→ 收入如何分配（9.3）→ 谁能分享数字红利（9.4）。
\end{examguide}

\section{技术进步与就业的经典理论}

第 8 章末尾的三个问题构成本章的展开顺序：技术进步替代了哪些人的工作（9.1 节）？就业结构发生了怎样的变化（9.2 节）？收入如何在劳动者之间重新分配（9.3 节）？又有哪些群体被挡在数字红利之外（9.4 节）？

\subsection{替代效应与创造效应}

技术进步对就业的影响沿两条渠道展开。\textbf{替代效应}：技术使机器完成原本由劳动承担的任务，劳动需求下降，19 世纪卢德分子捣毁织布机的恐慌，即对替代效应的朴素反应。这一恐慌在经济学上被称为\textbf{卢德谬误}（Luddite fallacy）：只看到替代效应而忽视创造效应，误将技术进步等同于就业总量下降；两百年来就业总量并未随自动化萎缩，技术进步改变的是就业结构而非总量。\mn{考试表述：技术进步对就业总量的净效应取决于替代与创造的相对强度。}\textbf{创造效应}：技术进步降低生产成本、提高真实收入、创造新产品与新任务，劳动需求上升。

\begin{derivation}{替代与创造的分解}{}
\textbf{替代与创造的局部均衡分解。}考虑技术进步（$A$ 上升）对劳动需求 $L$ 的影响。在局部均衡框架中，$A$ 上升使生产函数外移，其效应分解为：\textbf{价格效应}：成本下降使产品价格下降、需求扩大，劳动需求上升（创造渠道）；\textbf{替代效应}：给定产出水平，技术使单位产出的劳动投入下降（替代渠道）。形式化地，
\[
\frac{d\ln L}{d\ln A} = \underbrace{\varepsilon_D \cdot \gamma}_{\text{创造效应 } > 0} - \underbrace{(1 - \alpha)}_{\text{替代效应 } < 0},
\]
其中 $\varepsilon_D$ 为产品需求弹性，$\gamma$ 为技术进步向价格下降的传导率。净效应的符号取决于需求弹性与成本节约的相对大小：\important{需求弹性越大（价格下降撬动的需求扩张越强），创造效应越可能占优}。
\end{derivation}

这一分解的经验含义：需求弹性大的行业（价格敏感、市场广阔的产品）技术进步往往扩大就业，需求弹性小的行业则相反。数字技术的特点在于同时放大两端，自动化能力强化替代效应，网络外部性与新产品创造强化需求扩张，因此数字技术对就业总量的净效应不存在先验结论，这正是"AI 会否消灭就业"的争论长期无解的结构性原因。

替代恐慌在历史上反复上演：1811 年至 1816 年英国卢德分子捣毁织布机，1930 年代凯恩斯提出"技术性失业"的忧虑，1960 年代美国经历自动化恐慌，再到当前的 AI 失业论，每轮技术浪潮都伴随同构的恐惧。而两百年的经验事实是：就业总量未随自动化萎缩，就业结构却持续迁移，农业就业占比从 19 世纪的过半降至今日个位数，制造业随后收缩，服务业与知识劳动扩张。结构迁移本身构成调整成本：被替代岗位的劳动者需要转岗与再培训，"就业总量无虞"与"个体失业风险真实"并行不悖，这正是讨论技术冲击时的双重出发点。

\subsection{技能偏向型技术进步}

技术进步并非中性地影响所有劳动者：1970 年代以来，美国高技能与低技能劳动者的工资差距持续扩大，这一事实推动劳动经济学家追问——技术进步让谁受益？技能偏向型技术进步（SBTC）假说给出的回答是：技术本身是偏向性的，它系统性地提高了对高技能劳动的需求。

\begin{definition}[技能偏向型技术进步（SBTC）]
技能偏向型技术进步是指系统性地提高对高技能劳动相对需求的技术进步，在相同要素供给下，技术变化使高技能劳动的边际产出相对上升。
\end{definition}

\begin{derivation}{SBTC 与技能溢价}{}
\textbf{SBTC 与技能溢价的推导。}设生产函数为高技能劳动 $H$ 与低技能劳动 $L$ 的 CES 组合：
\[
Y = \bigl[(A_H H)^{\frac{\sigma-1}{\sigma}} + (A_L L)^{\frac{\sigma-1}{\sigma}}\bigr]^{\frac{\sigma}{\sigma-1}},
\]
其中 $A_H$、$A_L$ 为两类劳动的技术效率，$\sigma$ 为替代弹性。竞争性要素市场下，两类劳动的工资比等于其边际产出比：
\[
\frac{w_H}{w_L} = \frac{MP_H}{MP_L} = \left(\frac{A_H}{A_L}\right)^{\frac{\sigma-1}{\sigma}} \left(\frac{H}{L}\right)^{-\frac{1}{\sigma}}.
\]
取对数：
\[
\ln\frac{w_H}{w_L} = \frac{\sigma-1}{\sigma}\ln\frac{A_H}{A_L} - \frac{1}{\sigma}\ln\frac{H}{L}.
\]
即技能溢价由技术偏向（$A_H/A_L$）与相对供给（$H/L$）共同决定。
\end{derivation}

技能溢价的解读分两个维度。\textbf{技术维度}：当 $\sigma > 1$（两类劳动替代性较强）时，$A_H/A_L$ 上升（技术偏向高技能）提高技能溢价——\important{技能偏向型技术进步扩大高技能与低技能的工资差距}；当 $\sigma < 1$（互补性强）时方向相反。实证共识 $\sigma > 1$，SBTC 因而被确立为 1980 年代以来技能溢价上升的主要解释。\textbf{供给维度}：$H/L$ 上升（高技能供给增加）压低技能溢价，教育扩张是抑制工资差距的市场力量。\mn{技能溢价公式速记：$\ln\frac{w_H}{w_L}=\frac{\sigma-1}{\sigma}\ln\frac{A_H}{A_L}-\frac{1}{\sigma}\ln\frac{H}{L}$，前项技术偏向、后项供给抑制；符号判断先看 $\sigma$ 与 1 的大小。}技能溢价的长期走势取决于"技术偏向"与"教育供给"的赛跑：数字经济时代前者加速，后者相对滞后，这正是"大学生贬值与技能溢价并存"悖论的解释框架。

SBTC 的实证证据集中于教育溢价的长期上升：美国大学与高中学历劳动者的工资差由 1979 年的约 40\% 升至 1990 年代的约 60\% 并维持高位（Autor 与 Katz 等）。中国的对照经验更为复杂：高等教育扩招使高技能供给快速增加（$H/L$ 上升，压低溢价），但数字岗位的技能溢价仍然居高，"大学生贬值与技能溢价并存"恰是技术偏向与教育供给赛跑的中国版本。SBTC 的局限也在此显现：它刻画要素层面的偏向，无法解释为何被替代的恰是中等技能岗位而非最低技能岗位，任务模型（9.1.3 节）补上了这一层。

\subsection{任务模型}

Acemoglu 与 Restrepo 的任务模型将技术进步刻画为\textbf{自动化边界的移动}：生产由连续统任务构成，每项任务可由劳动或资本（自动化）完成。任务的连续统结构意味着技术进步不再"整体替代劳动"，而是"逐任务蚕食"，就业影响取决于自动化边界移动的方向与速度。

\begin{derivation}{自动化边界}{}
\textbf{自动化边界的比较静态。}设任务区间为 $[0,1]$，任务 $i$ 的产出为 $x(i)$。自动化边界为 $\bar\theta \in (0,1)$：任务 $i \leq \bar\theta$ 已被自动化（资本完成），$i > \bar\theta$ 仍由劳动完成。劳动需求为
\[
L^d = \int_{\bar\theta}^{1} \ell(i)\, di,
\]
其中 $\ell(i)$ 为任务 $i$ 的劳动投入。自动化边界外移（$\bar\theta$ 上升）的直接效应：劳动需求减少（替代效应）。但存在两个抵消渠道：\textbf{生产率效应}：自动化降低成本，产出扩大，未自动化任务的劳动需求上升；\textbf{新任务创造}：技术同时创造新任务（区间右端延伸），吸收被替代的劳动。模型的核心命题：\important{就业的长期走势取决于"自动化任务"与"新任务创造"的竞赛}。
\end{derivation}

数字时代任务模型的要义在 AI 处变得锋利：历史上自动化局限于体力与程序性任务，而 AI 的自动化能力覆盖认知任务，替代效应首次系统性地指向白领与知识劳动，第 14 章对 AI 冲击的讨论即建立在本节框架之上。任务模型同时给出评估自动化的判据：关注的重心不应是"机器能做多少"，而是"新任务创造的速度能否追上自动化边界的外移"——生产率效应与新任务创造是就业的两条出路。

任务模型的预测在中国制造业得到直观印证：中国自 2013 年起保持全球最大工业机器人市场，2023 年工业机器人安装量约占全球一半以上（国际机器人联合会），自动化边界在装配、焊接等程序性任务上快速外移；与此同时，数据分析、算法标注、直播运营等新任务持续创造就业。\mn{速查锚点：中国自 2013 年起保持全球最大工业机器人市场，2023 年安装量占全球一半以上（国际机器人联合会），"机器换人"论述题的标准开场数据。}"机器换人"与"新职业涌现"在中国同期发生，正是自动化与新任务创造竞赛的微观样本。任务模型的政策含义：与其争论自动化的方向，不如降低任务转换成本（技能培训、再就业服务），让劳动者能够跨越被自动化的任务区间。

\section{数字经济与就业结构}

\subsection{就业极化}

数字技术对就业结构的显著效应是就业极化。1990 年代以来，美国等发达经济体的就业份额呈现"两头大、中间小"的演进：高技能与低技能岗位扩张，中等技能的程序性岗位萎缩。任务模型（9.1 节）为此提供了可检验的机制，本节将其形式化。

\begin{definition}[就业极化（job polarization）]
就业极化是指高技能（抽象、创造类）与低技能（面对面服务类）就业扩张、中间技能（程序性、规则化任务）就业萎缩的就业结构变化：就业分布呈现\important{"两头大、中间小"}的形态。
\end{definition}

\begin{derivation}{就业极化的任务分布解释}{}
\textbf{极化的替代弹性差异。}将任务按性质分为三类：抽象任务（$A$，管理、创造）、程序任务（$R$，记账、装配）、手工任务（$M$，清洁、护理）。设每类任务由劳动与资本（自动化）以 CES 技术组合：
\[
Y_j = \bigl[(A_j L_j)^{\rho_j} + (B_j K_j)^{\rho_j}\bigr]^{\frac{1}{\rho_j}}, \qquad j \in \{A,\, R,\, M\},
\]
资本—劳动替代弹性 $\sigma_j = 1/(1-\rho_j)$。程序任务规则化程度最高、最易编码，替代弹性最大（$\sigma_R > 1$）；抽象任务需要判断与创造、手工任务需要物理灵活性与情境适应，均难以自动化（$\sigma_A < 1$，$\sigma_M < 1$）。令 $z = B_j r / (A_j w)$ 为资本相对劳动的有效成本比，劳动成本份额
\[
s_j^L = \frac{w L_j}{w L_j + r K_j} = \frac{1}{1 + z^{\,\sigma_j - 1}}.
\]
自动化使资本效率上升（$B_j$ 上升，$z$ 上升），劳动份额的响应为
\[
\frac{\partial \ln s_j^L}{\partial \ln z} = -(\sigma_j - 1)\, s_j^K,
\]
$\sigma_j > 1$ 时劳动份额下降（且随 $\sigma_j - 1$ 放大），$\sigma_j < 1$ 时反而上升。\important{极化条件：$\sigma_R > 1 > \sigma_A, \sigma_M$}，自动化冲击下程序任务劳动需求收缩、抽象与手工任务劳动需求扩张，就业分布呈 U 形。
\end{derivation}

替代弹性差异的传导机制：自动化边界（9.1 节）首先跨越程序任务区，中间技能的劳动需求系统性下降。数据上，1980 年代以来美国就业在工资分布两端扩张、中间收缩的 U 形模式，与这一理论预测高度吻合。

极化的跨国证据高度一致：美国就业份额在工资分布两端扩张、中间收缩的 U 形模式自 1980 年代持续（Autor 与 Dorn），欧洲劳动力市场同样呈现例行任务就业的相对萎缩（Goos 与 Manning），计算机化对"例行任务"的替代是共同的解释变量。中国的图景存在差异：制造业转移与城镇化吸收了大量中低技能劳动力，极化的形态存在但被缓冲（吕世斌与张世伟，2015；李宏兵等，2017），程序性岗位（会计、银行柜员、客服）的收缩与数字服务岗位的扩张已经显现，且叠加了区域差异（9.4.3 节）。极化的福利含义：被挤出的中等技能劳动者向手工任务迁移，不仅压低自身工资，还挤占低技能劳动者的就业空间，中间塌陷的代价由两端共同承受。

\begin{figure}[htbp]
\centering
\begin{tikzpicture}[>=Stealth, scale=0.95]
\begin{axis}[
    width=10.5cm, height=5.6cm,
    xlabel={技能水平（低 $\rightarrow$ 高）}, ylabel={就业变化率},
    xmin=0, xmax=10, ymin=-0.25, ymax=0.2,
    axis lines=left,
    xtick={1.5,5,8.5}, xticklabels={低技能（手工任务）, 中等技能（程序任务）, 高技能（抽象任务）},
    x tick label style={font=\small, align=center, text width=2.4cm},
    yticklabels={},
]
\addplot[color=main, very thick, domain=0:10, samples=200] {0.25*(x-5)^2/25 - 0.15};
\draw[dashed, gray] (axis cs:5,0.2) -- (axis cs:5,-0.25);
\end{axis}
\end{tikzpicture}
\caption{就业极化}
\label{fig:polarization}
\end{figure}

图~\ref{fig:polarization} 以 U 形曲线刻画就业极化的经验形态：中间技能（程序任务）的就业萎缩，两侧（手工任务与抽象任务）扩张。值得注意的是极化与技能溢价的组合效应：中等技能劳动者向下迁移（挤入手工任务）与向上受阻（抽象任务有技能门槛）同时发生，就业结构的变化由此传导为收入分布的变化——这正是 9.3 节分配分析的起点。

就业极化并不必然伴随工资极化，这一辨析是简答与论述题的常见区分点。理论上存在两股方向相反的力量：其一，\textbf{消费互补性}，高技能集聚城市的收入上升拉动对低技能服务的需求，抬高手工任务工资（Autor 与 Dorn，2013）；其二，\textbf{供给转移}，被程序任务挤出的中等技能劳动者涌入手工任务，压低低技能工资。两股力量的净效果不确定，因此\important{就业极化是实证上稳健的现象，工资极化在理论与证据上都弱得多}：美国 1980 年代以来工资分布的低端并未出现与就业对称的扩张（Autor、Katz 与 Kearney，2006）。Acemoglu 与 Restrepo（2017）进一步区分了短期与长期：自动化在短期内压低受替代任务密集部门的工资，长期则经由新任务创造部分回补，工资极化的证据强度取决于观察窗口。辨析题的作答结构由此固定：就业结构看任务替代弹性（$\sigma_R > 1 > \sigma_A, \sigma_M$），工资结构看消费互补与供给转移的净效果。

\subsection{技能需求逆转}

极化描述的是技能需求在结构上的移动方向，而 2000 年以来的证据提出了一个更尖锐的命题：对认知技能的需求本身可能放缓甚至逆转。这一命题动摇了"教育与技术的赛跑"框架的默认前提，是学硕热点的辨析素材。

\begin{definition}[技能需求逆转（skill demand reversal）]
技能需求逆转是指 2000 年前后发达国家对认知技能与高等教育技能的需求增长放缓乃至停滞：高技能劳动者沿职业阶梯下移、挤出低技能者，认知任务的劳动报酬下降，大学溢价增长放缓。
\end{definition}

Beaudry、Green 与 Sand（2016）的假说：计算机化浪潮在 1990 年代到达顶峰后，其组织变革（组织结构重构）告一段落，高技能需求的结构性推动力消退，高技能劳动者向下挤出低技能者，形成"繁荣-萧条"式调整。经验证据呈现三个特征：其一，就业率与平均工资的关系从正相关转为"逆 C"形，中等学历就业者被上下两端挤压；其二，美国认知任务的实际报酬下降约 2\%、手工任务下降约 8\%，高等教育的工资溢价仍在扩大但增速显著放缓（研究生溢价与大学溢价出现分化）；其三，2000 年后美国大学生溢价增长停滞，而研究生溢价继续上升。\mn{技能需求逆转证据三件套：逆 C 形态、认知任务报酬下降、大学溢价增速放缓（研究生溢价继续上升）。文献锚点：Beaudry–Green–Sand（2016）。}

逆转的三条机制各有文献支撑。其一，\textbf{技术投资放缓}：Beaudry–Green–Sand 以组织资本模型刻画技术繁荣的退潮，技能需求随组织重构的完成而回落；其二，\textbf{软件创新}：Aum（2017）与 Aum 和 Shin（2020）指出软件的模块化与标准化降低了对高技能编程与工程技能的需求；其三，\textbf{高技能自动化}：Acemoglu 与 Restrepo（2018）的框架把任务模型从低技能自动化推广到认知任务的自动化，技能需求的方向取决于自动化与新任务创造的相对速度。三条机制分别对应技术生命周期、技术形态与任务结构三个层次，第 14 章 AI 对认知任务的切入正是第三条机制的当代延伸。

\begin{misconception}
\textbf{误区}："技能需求逆转等于读书无用"。逆转的准确含义是对认知技能的需求\textbf{增速}放缓，而非技能溢价消失：大学溢价在多数国家仍然为正，研究生溢价甚至继续上升。逆转改变的是技能分布内部的相对需求结构，教育的私人回报与信号价值依然成立；对政策而言，启示不是减少教育，而是教育内容从"更多年限"转向"技能组合的质量"。
\end{misconception}

中国的证据尚在形成之中。基于中国劳动力市场数据的测算显示，1990 年代末起就业极化开始出现、2008 年前后趋于明显，2010 年前后高技术行业出现需求放缓的早期迹象（吕世斌与张世伟，2015；李宏兵等，2017），信息传输与软件业的高学历比重在 2010 年前后的波动被部分文献解读为逆转信号。需要说明的是，中国同时经历着制造业转移、城镇化与数字化的三重叠加，识别出"纯粹的技能需求逆转"远比美国困难，这一结论在考试中表述为"有迹象、有缓冲、待检验"即可。技能需求逆转与极化构成完整图景：\important{极化刻画需求结构的形状，逆转刻画需求结构的位移，两者共同决定技能溢价与就业结构的中期走势}。

\subsection{零工经济与平台用工}

平台将劳动的组织形式从"雇员—雇主"推向"接单者—平台"的零工经济。零工岗位（骑手、网约车司机、内容创作者）已构成数字平台用工的主体，其供给决策与保障缺失构成就业分析的制度焦点。

\begin{definition}[零工经济（gig economy）]
零工经济是指劳动者以按单计酬、任务制方式与平台合作、而非建立雇佣关系的劳动形态：骑手、网约车司机、内容创作者是其典型群体。
\end{definition}

\begin{derivation}{零工劳动供给}{}
\textbf{零工劳动供给的权衡模型。}设劳动者选择正规就业或零工就业。正规就业的效用为 $U_E = w_E + b$，其中 $w_E$ 为工资、$b$ 为保障价值（社保、稳定预期）；零工就业的效用为 $U_G = w_G(\text{工作时长}) + f(\text{灵活性})$，其中 $w_G$ 为按单收入、$f$ 为灵活性价值（自主安排时间）。劳动者选择\important{零工}的条件为
\[
w_G + f \geq w_E + b.
\]
即零工的总效用（按单收入 + 灵活性价值）不低于正规就业（工资 + 保障价值）时，零工供给发生。
\end{derivation}

零工经济的吸引力在于 $f$（灵活）与进入门槛低；其代价是 $b$ 的缺失：平台以"承揽"而非"雇佣"定性与劳动者合作，社保、最低工资、解雇保护的雇主责任被规避。这一模型的制度焦点在于 $b$ 的定价：当平台不承担雇主责任时，$b$ 被外部化给劳动者与社会（职业伤害、养老缺口的社会化），零工定价因此产生负外部性，这是平台用工关系争议（雇佣 vs 承揽）的经济学核心。中国《关于维护新就业形态劳动者劳动保障权益的指导意见》（2021）引入"不完全符合确立劳动关系情形"的第三类身份，正是对 $b$ 部分补偿的制度尝试。\mn{平台用工三种关系形态：劳动关系（骑手 vs 直营）、不完全劳动关系（网约车司机 vs 平台）、民事承揽关系（自由接单者）。2021 年指导意见创设的"不完全劳动关系"是制度创新考点。}三种形态的判定以劳动从属性（人身、经济、组织从属性）为尺度，"不完全劳动关系"居于两者之间，这一制度空间正是平台用工争议的落点。

零工经济的规模已构成就业的基本盘。人社部口径的灵活就业人员约 2 亿，全国总工会统计的新就业形态劳动者约 8400 万（2023），仅美团平台年获得收入的骑手即超七百万，网约车司机达数百万量级。零工平台的中国实践还展现了算法的深度嵌入：派单算法决定接单机会，路线算法压缩配送时间，评价算法沉淀服务质量，平台通过数据与算法对劳动者实施"准管理"，却以承揽关系规避雇主责任，劳动者实质上处于"形式独立、实质从属"的状态，这正是劳动关系认定的难点所在。

制度回应在劳动保护与算法治理两个方向推进：2021 年指导意见确立第三类身份后，2022 年起北京、上海、江苏、广东等 7 省市试点新就业形态就业人员职业伤害保障（覆盖骑手、网约车司机等），外卖平台在监管推动下公开配送算法、设定最长在线时长，对 $b$ 的补偿从"身份承认"走向"待遇落地"。经济学的判断标准始终如一：当平台行使实质性管理权力时，责任不应被合同形式外包。就业一侧的分析至此完整，结构变化（极化）与用工形态变化（零工）共同作用于劳动者的收入形成，9.3 节转入收入分配的三个机制。

\section{数字经济与收入分配}

\subsection{超级明星效应}

数字经济在收入端的标志性现象是\textbf{超级明星效应}：市场奖励向少数顶尖供给者极度集中，头部主播的收入超过行业总量的绝大部分，头部程序员与普通程序员的差距远超传统行业。

\begin{derivation}{超级明星模型}{}
\textbf{Rosen 超级明星模型的核心直觉。}设市场对质量为 $q$ 的供给者的支付意愿为 $p(q)$，供给者向 $n$ 个消费者提供服务的边际分发成本为 $c$（传统市场为正，数字市场趋近于零）。质量为 $q$ 的供给者的收入为
\[
R(q) = p(q) \cdot n - c \cdot n.
\]
传统市场中，个人服务的市场规模受物理约束（一场演出、一次面诊），$n$ 有上限，质量的收入弹性有限。数字技术使 $n$ 急剧扩张（直播、课程、代码的分发边际成本趋于零），收入对质量的弹性被放大：
\[
\frac{dR}{R} = \frac{dp}{p} + \frac{dn}{n},
\]
质量微小优势转化为市场规模的指数级差距。当质量优势与网络外部性结合（粉丝效应、推荐算法放大），收入分布收敛于\important{幂律}：\important{超级明星获得"赢家通吃"的收入，尾部供给者的收入被压缩}。
\end{derivation}

这一机制与第 3 章的网络外部性同构：超级明星效应是需求方规模经济在劳动市场的对应物，平台的网络效应与算法的推荐机制放大质量差异，使头部供给者获得不成比例的需求。超级明星效应的分配含义：收入分布呈幂律而非正态，中位数收入远低于均值，"1\% 的人拿走大部分"在数字劳动市场成为常态；同时它对供给者形成争夺"头部"位置的投入竞争，尾部收益在竞争中进一步恶化。

超级明星效应的中国样本集中在直播电商：头部主播单场 GMV 可达数亿元，超过多数零售企业全年的线上销售额，而中长尾主播的收入接近普通服务业水平，带货集中度与平台的双边抽成结构（第 5 章）相互强化。对照音乐产业的全球格局（头部艺人占据流媒体播放量的绝大部分），数字分发成本趋零使"赢家通吃"从娱乐业扩散到一切可数字化服务的行业。政策层面的关注点有二：其一，超级明星收入与劳动投入的脱钩是否构成激励扭曲（对头部的过度投入竞争）；其二，幂律分布下收入统计的误导性（均值与中位数严重背离），后者直接进入收入度量的讨论。

\subsection{劳动收入份额的下降}

超级明星效应的宏观对应是\textbf{劳动收入份额}（labor share）的下降：增加值中分配给劳动的比例在多数发达经济体持续走低。Karabarbounis 与 Neiman（2014）对全球样本的测算显示，劳动收入份额由 1970 年代的约 66\% 降至 2010 年代的约 59\%，这一趋势构成收入分配研究的重要经验事实。

\begin{derivation}{劳动收入份额}{}
\textbf{加成率与劳动收入份额的代数关系。}设总产出 $Y$，增加值中劳动收入为 $WL$。劳动收入份额为
\[
\lambda = \frac{WL}{PY}.
\]
在竞争性均衡（$P = MC$）下，$\lambda$ 由技术参数（生产函数弹性）决定；当市场势力使价格高于边际成本（$P > MC$，加成率 $\mu = P/MC > 1$）时，垄断利润挤占要素收入。Autor 等的"超级明星企业"假说的代数表达：若加成率上升与销售额向高加成企业集中同时发生，则
\[
\lambda = \sum_i \lambda_i s_i \quad \text{随 } s_i \text{向高加成企业集中而下降},
\]
其中 $s_i$ 为企业 $i$ 的销售份额。
\end{derivation}

\important{劳动收入份额下降不必来自"资本欺负劳动"的传统机制，而是来自产品市场竞争弱化：高加成企业（平台、科技巨头）的市场份额扩大，即使其内部的 $\lambda_i$ 不变，总体的 $\lambda$ 也下降。}数字经济的相关性在于：网络外部性与数据壁垒（第 3、6 章）正是高加成企业市场份额扩张的放大器，超级明星企业机制与超级明星效应（9.3.1 节）在劳动市场上互为镜像。

中国劳动报酬份额的走势提供了另一视角：国家统计局口径下，劳动报酬占 GDP 比重由 2000 年前后的约 50\% 降至 2007 年前后的低点，此后回升并保持相对稳定，低于多数发达经济体同期水平。\mn{速查锚点：全球劳动收入份额自 1980 年代初约 65\% 降至 2000 年代末 60\% 以下（Karabarbounis 与 Neiman, 2014）。}下降段与出口导向工业化、资本密集化同步，回升段与劳动力供给拐点重合，说明劳动份额由要素相对供给与技术结构共同决定，数字平台的扩张则通过高加成率渠道（上文）对其构成新压力。劳动份额的度量本身也有争议：个体经营者收入（网约车司机、小商户）在统计中计入营业盈余还是劳动报酬，直接改变份额读数，"劳动者—平台"新型关系对国民账户分类的冲击，是劳动份额研究在数字经济时代的新课题。

\subsection{数据租金与分配}

数据要素进入生产函数（第 8 章三要素核算）后，一个分配问题随之而来：数据产生的超额收益归谁？数据由劳动者与用户的行为产生（每一次点击、每一单配送都是数据投入），但数据资产通常由平台持有（第 7 章三权分置的制度安排），租金的归属由此成为劳资分配的新焦点。

\begin{derivation}{数据租金的分割}{}
\textbf{数据租金在资本与劳动间的分割。}沿用第 8 章的三要素生产函数，并设规模报酬不变：
\[
Y = A K^\alpha L^\beta D^\gamma, \qquad \alpha + \beta + \gamma = 1,
\]
竞争性分配下各要素按其产出弹性取得收入：$wL = \beta Y$、$rK = \alpha Y$、数据租金 $r_D D = \gamma Y$。数据租金归谁取决于数据资产的所有权配置：若数据由劳动者与用户产生但归平台（资本）持有，则资本实际所得为
\[
\frac{rK + r_D D}{Y} = \alpha + \gamma, \qquad \frac{wL}{Y} = \beta = 1 - \alpha - \gamma.
\]
即\important{数据要素在总产出中的份额 $\gamma$ 越大、且数据租金归资本时，劳动收入份额越低}。
\end{derivation}

数据租金的分割把 9.3.2 节的劳动份额下降补充为两个渠道的合成：市场势力渠道（加成率上升）与要素构成渠道（数据租金归资本）。\mn{数据租金速记：资本实际份额 $\alpha+\gamma$、劳动只剩 $\beta = 1-\alpha-\gamma$；记法：谁持有数据资产，谁就拿走 $\gamma$ 这块租金。}后者的特殊之处在于数字经济中"数据由谁产生"与"数据归谁所有"的分离：劳动者与用户贡献了数据这一投入品，却因数据资产被平台独占而无从分享其租金，这与第 7 章三权分置"淡化所有权、强化使用权"的制度安排形成张力。数据要素收益分配制度的改革方向（"谁投入、谁贡献、谁受益"的初次分配原则）正是在 $\gamma$ 的归属上做文章：承认劳动对数据形成的贡献，让数据收益的一部分按贡献回归劳动者。

数据要素收益分配的制度探索正在展开。"数据二十条"提出建立"谁投入、谁贡献、谁受益"的初次分配机制，强调数据要素收益向数据价值创造者倾斜；2024 年企业数据资源会计处理新规实施，数据资产入表使数据收益显性化，为租金分割提供了计量基础。但现实中的分割仍然失衡：平台基于用户行为数据训练推荐算法并获取广告与佣金收入，骑手轨迹数据被用于优化派单效率，劳动者与用户既未获得数据收益的分配，也缺乏对数据使用的议价权，数据租金归资本的格局尚未被制度打破，这是第 7 章数据产权讨论在分配维度的延伸。

\subsection{平台抽成与劳动者分配}

平台对零工劳动者的抽成（佣金率）构成分配分析的另一焦点。佣金率的经济学（第 5 章）与竞争瓶颈（第 4 章）在此叠加：劳动者侧的\textbf{多归属成本}（同时注册多个平台的时间成本、账号管理成本）使其接近单归属状态，而劳动者对平台派单与流量的依赖远超传统分包关系，越过第 4 章竞争瓶颈的判断线，价格结构翻转为平台向劳动者侧收取高于垄断加成的费用；同时算法掌握的需求分配权进一步加深了这种依赖。其分配后果可以形式化。

\begin{derivation}{佣金归宿与劳动净收入}{}
设客单价为 $p$，劳动者完成的单量为 $q$，平台佣金率为 $t$，劳动者的装备与运营成本为 $c(q)$（$c' > 0$）。劳动者的净收入为
\[
w_{\text{净}} = (1 - t)\, p\, q - c(q).
\]
佣金上调 $\Delta t$ 后，净收入变化由两部分构成：给定单量下的佣金支出增加 $-p q\,\Delta t$，以及单量收缩带来的收入损失 $(1-t)p\,\mathrm{d}q$。由单量供给弹性 $\varepsilon_{q,t} = -\frac{\mathrm{d}q/q}{\mathrm{d}t/t} > 0$ 得 $\mathrm{d}q/\mathrm{d}t = -\varepsilon_{q,t}\, q/t$，代入得
\[
\frac{\partial w_{\text{净}}}{\partial t} = -p\, q \cdot \left[1 + \underbrace{(1-t)\,\frac{\varepsilon_{q,t}}{t}}_{\text{单量收缩的放大}}\right],
\]
第一项为佣金的直接负担，第二项为单量收缩的间接损失。两个极端情形界定佣金的实际归宿：\textbf{完全无弹性}（$\varepsilon_{q,t} = 0$，劳动者别无选择）时，佣金上调一单位，净收入按客单价 × 单量全额下降，佣金全部由劳动承担；\textbf{高弹性}（$\varepsilon_{q,t}$ 大，劳动者可转投他业）时，单量收缩更剧烈，平台为维持司机供给被迫提高客单价或下调佣金率，部分负担在均衡中转嫁给消费者与平台，但转嫁以劳动者的退出威胁为前提。
\end{derivation}

零工劳动者单归属结构（多归属成本高）的含义是 $\varepsilon_{q,t}$ 被锁定在低位，佣金的分配归宿高度偏向劳动者承担。"算法剥削"争论的经济学实质由此清晰：算法掌握的需求分配权压低劳动者对单量的议价能力，单归属结构压低其退出威胁，二者共同使 $\varepsilon_{q,t} \to 0$，佣金成为劳动者的刚性负担。

抽成透明度已成为监管焦点：2021 年起交通运输部等部门要求网约车平台公开计价规则与抽成比例，主要平台相继上线"司机收入透明化"功能；外卖平台在 2021 年外卖员权益保障新政后下调高峰时段抽成、设立商家与骑手的协商机制。这些措施的经济学意义在于把佣金议价从"信息黑箱"推向"可谈判"：透明度降低劳动者的信息不对称（单归属状态下平台独占定价信息），协商机制则在事实上承认劳动者的部分退出威胁。不过，只要多归属成本与算法分配权不变，佣金归宿的总体格局就难以根本改变。

\section{数字鸿沟}

\subsection{三层结构}

数字鸿沟回答的是收入分配问题的最后一环：即使数字红利在总量上扩大、在劳动者之间被分配，仍有群体根本未被纳入数字经济的收益范围。第 8 章增长框架中的数字鸿沟压低落后地区的稳态（8.4.2 节），本节剖析其微观结构。

\begin{definition}[数字鸿沟]
数字鸿沟是指不同群体在数字技术获取、使用与收益转化上的\important{系统性差距}：接入数字技术的机会、利用数字技术的能力、以及数字红利向生活机会的转化，在不同群体之间呈现结构性分化。
\end{definition}

其结构分为三层。\textbf{接入沟}：能否接入互联网与数字设备（基础设施层面），随着覆盖率提升，接入沟在多数地区趋于弥合。\textbf{使用沟}：接入之后能否有效利用（技能层面）：同一互联网面前，有人用于学习与创业，有人仅用于娱乐；使用沟的分化程度与教育、收入高度相关。\textbf{结果沟}：数字利用差异转化为生活机会的差异（收益层面）：教育、健康、收入、参政能力的系统性差距，是数字鸿沟的最终形态。\mn{数字鸿沟三层记忆：接入（能不能上）、使用（会不会用）、结果（有没有受益）。政策含义：接入沟靠基建（普遍服务），使用沟靠教育（数字素养），结果沟靠制度（包容性应用设计）。}三层之间的递进与转化关系（9.4.2 节）构成鸿沟治理的政策次序。

\begin{derivation}{数字接入决策}{}
\textbf{数字接入的投资决策模型。}设个体接入数字技术的收益为 $B(\theta, G)$（$\theta$ 为个体技能，$G$ 为本地基础设施质量），接入成本为 $c$。个体接入的条件为
\[
B(\theta, G) \geq c.
\]
收益不低于接入成本时接入发生。
\end{derivation}

政策的两个杠杆由此读出：降低 $c$（补贴终端与资费，即接入沟政策）与提高 $B$（技能培训提升 $\theta$、数字公共服务提升 $G$，即使用沟政策）。模型的含义是：\important{单靠接入补贴无法弥合鸿沟}，若 $B(\theta, G)$ 因技能不足而接近于零，则无论 $c$ 降到多低，接入也不发生或不可持续；接入之后的使用与结果差距，才是鸿沟的主体。

接入沟的全球图景：国际电信联盟估计（2023）全球约三分之二人口使用互联网，最不发达国家仅约三分之一；中国农村互联网普及率约 66\%（CNNIC，2023），与城镇的约 83\% 相差约十七个百分点，但差距持续收窄，接入沟处于"趋于弥合"阶段。值得注意的是老龄鸿沟：老年人因数字技能不足而被公共服务线上化排斥（2020 年疫情期间健康码使用困难的典型事件，国务院随即出台老年人运用智能技术的实施方案），同一接入之下使用能力的分化使使用沟在年龄维度最为尖锐。使用沟与结果沟的度量也日益精细，数字素养测评、数字生活参与指数等陆续进入统计体系，为鸿沟治理提供靶点。

\subsection{三层结构的动态关系}

三层结构之间存在动态递进关系，这是数字鸿沟分析中最容易被忽略的要点。接入沟的弥合不等于鸿沟的消失：当"能不能接入"的差距缩小后，"会不会使用"的差距成为鸿沟的主体；而当使用差距缩小时，结果差距（数字红利转化为生活机会的能力）继续分化。\important{鸿沟的形态随数字化的深入不断升级，而非被技术进步自动消除}，每一层鸿沟的弥合都会暴露下一层更深的分化，这是数字鸿沟研究的基本事实。

\begin{derivation}{数字鸿沟的自我强化}{}
数字鸿沟的持续性来自其自我强化机制。设个体的数字技能 $s_t$ 通过使用数字技术积累：$s_{t+1} = s_t + \mu(s_t)$，其中 $\mu'(s_t) > 0$，技能越高，从使用中学习的边际收益越大（"用得多才学得快"）。数字技术的收益 $B(s_t)$ 随技能凸性增长（高技能者从数字服务中提取的价值不成比例地高：高技能者用网络学习与创业，低技能者用于娱乐消费）。于是技能差距按
\[
\frac{d(s_H - s_L)}{dt} = \mu(s_H) - \mu(s_L) > 0
\]
持续扩大，初始的技能差距通过使用行为的差异自我放大——数字鸿沟的本质是\important{初始不平等经由数字技术的收益凸性被放大}。
\end{derivation}

这一机制的推论：干预的时点越早越好（在差距被放大之前介入），且单纯接入（提高 $s$ 的使用机会）不改变 $\mu(\cdot)$ 的凸性，使用技能的训练才是切断自我强化链的关键。

动态递进的典型样本是 2020 年疫情期间的线上教育：接入层面城乡差距已不显著（通网率普遍较高），但"停课不停学"暴露出设备与使用差距（部分学生无智能手机或需与家人共享、家长无力辅导），进而转化为学业成绩与升学机会的差距，接入弥合后暴露使用沟，使用沟又转化为结果沟，三层鸿沟在同一事件中完成一次完整的递进。数字鸿沟的代际传递同样值得警惕：父母的数字素养影响子女的数字接触（设备管理、学习监督），技能差距由此跨代累积，使"干预时点越早越好"的结论延伸到家庭教育层面。

\subsection{城乡与区域维度}

中国数字鸿沟的空间形态呈现城乡与区域双重分化：农村与中西部地区的数字基建、数字技能、数字红利均落后于城市与东部。农村互联网普及率长期低于城镇（约 66\% 对 83\%，CNNIC，2023），且差距在使用深度上更为显著：农村网民的信息获取、在线教育、数字政务使用率全面落后，娱乐类使用比例反而更高（使用沟的典型形态），农村电商创业者与城镇同行的差距不在设备而在运营能力。中西部与东部的差距则叠加了产业结构因素：数字经济的高薪就业与创业机会集中于服务业发达地区（杭州、深圳等数字产业集群），中西部即便完成接入，数字红利的本地转化渠道仍然有限，结果沟呈现空间固化。

数字鸿沟的测度已从单一接入率走向多维指数。国际电信联盟的 ICT 发展指数以接入、使用、技能三维度量各国差距；国内研究则常用城乡互联网普及率差、数字素养测评与数字生活参与指标刻画三层鸿沟。测度的意义在于政策的靶向：接入差距看基建投资回报，使用差距看教育投入，结果差距看包容性制度设计，三层鸿沟对应三类政策工具，测度因此是治理的起点。

政策应对构成组合拳，且各有其经济学定位。数字乡村建设（基建下沉）针对接入沟，其依据是第 8 章基础设施外部性论证。数字乡村战略从 2019 年《数字乡村发展战略纲要》到 2022 年《数字乡村发展行动计划（2022—2025 年）》持续加码；2015 年起实施的电信普遍服务机制以"中央财政补助 + 企业交叉补贴"的方式将光纤与 4G 覆盖推进到行政村，2021 年通光纤与通 4G 的行政村占比均超 99\%，接入沟的最后一公里主要靠公共投资完成。

平台经济下乡（电商进村、直播助农）则以商业模式而非补贴实现数字红利的空间扩散，其效率优势在于让市场激励而非行政指令完成资源动员：电商平台降低农产品上行门槛，"新农人"直播将县域特色产品接入全国市场，快递进村随之补位。其局限同样来自平台逻辑：下乡红利的分配仍受平台佣金结构约束（9.3.4 节），流量分配向头部集中（9.3.1 节），普通农户在直播助农中更多扮演供货者而非受益者。平台下乡缓解了接入与销售瓶颈，却未必改变红利的分布结构，这正是"数字红利扩散的条件性"在空间维度的体现。

数字素养教育针对使用沟，按自我强化模型（9.4.2 节）的逻辑，这是阻断鸿沟代际传递的关键环节：面向农民的电商培训、面向老年人的智能技术帮扶、面向农村学生的数字素养课程，本质都是在改变技能积累的凸性而非仅提供接入。治理的次序由此清晰：接入靠公共投资，使用靠教育培训，结果靠制度设计，三层鸿沟逐层对应，越往后的治理越难、周期越长，也越决定数字红利的最终归属。

第 8 章条件收敛的结论在此衔接：数字红利扩散以互补条件（物流、人力资本、制度）为前提，单纯"通网"不自动产生红利；同样，鸿沟治理若只停留在接入层面，也无法改变红利分配的分化趋势。城乡与区域数字鸿沟的治理，本质上是把第 8 章的互补条件与三层鸿沟的分层整合为政策次序：先补基础设施（接入），再育人力资本（使用），后调分配制度（结果）。

\begin{chapterbox}
\textbf{本章考点清单}

\begin{enumerate}
\item 替代效应与创造效应的分解（$\varepsilon_D \gamma$ 与 $1-\alpha$）\important{【专硕·核心】}
\item 卢德谬误的历史与理论辨析 \important{【专硕·核心】}
\item SBTC 模型：$\ln(w_H/w_L) = \frac{\sigma-1}{\sigma}\ln(A_H/A_L) - \frac{1}{\sigma}\ln(H/L)$ 的完整推导 \important{【学硕·热点】【专硕·核心】}
\item 任务模型：自动化边界、生产率效应与新任务创造 \important{【专硕·核心】}
\item 就业极化的任务分布解释与极化条件（$\sigma_R > 1 > \sigma_A, \sigma_M$）\important{【专硕·核心】}
\item 就业极化与工资极化的辨析（消费互补 vs 供给转移）\important{【专硕·核心】}
\item 技能需求逆转（Beaudry–Green–Sand，2016）：逆 C 形态、认知任务报酬下降与三条机制 \important{【学硕·热点】}
\item 零工劳动供给模型 $w_G + f \geq w_E + b$ 与"不完全劳动关系"制度 \important{【专硕·核心】}
\item Rosen 超级明星模型与幂律收入分布 \important{【专硕·核心】}
\item 劳动收入份额下降的"超级明星企业"机制 \important{【学硕·热点】}
\item 数据租金在资本与劳动间的分割（$\alpha + \gamma$ 与 $\beta$）\important{【专硕·扩展】}
\item 平台抽成对劳动者的分配效应 \important{【专硕·核心】}
\item 数字鸿沟三层结构与接入决策模型 \important{【专硕·核心】}
\end{enumerate}
\end{chapterbox}

就业与分配的分析反复印证同一个结论：数字红利的分享是条件性的——它要求技能与任务的匹配、制度对新型劳动关系的覆盖、以及对鸿沟的主动弥合。同样的"条件性"逻辑，在数字经济的金融维度将以另一种面貌出现：技术能降低金融服务的门槛，也能放大金融体系的固有风险。
